%META {"question":"Trên đường thẳng a, hai điểm A và B cùng nằm một phía của điểm O. Tia OA và tia OB là hai tia gì?","options":["Trùng nhau — nên không phải hai tia đối nhau","Đối nhau","Không chung gốc O","Vuông góc với nhau"],"answer":0,"note":"A và B cùng phía của O nên tia OA và tia OB cùng hướng, tức trùng nhau. Hai tia đối nhau phải nằm khác phía của gốc chung."}
\documentclass[border=5pt]{standalone}
\usepackage{tkz-euclide}
% Màu dùng chung cho toàn bộ hình vẽ (ảnh stack với nền tối của web)
\definecolor{tminc}{RGB}{226,232,240}   % chữ + nét chính
\definecolor{tmblue}{RGB}{0,210,255}    % nhấn neon xanh
\definecolor{tmpink}{RGB}{255,0,200}    % nhấn neon hồng
\definecolor{tmgreen}{RGB}{0,255,136}   % nhấn neon xanh lá
\color{tminc}

\begin{document}
\begin{tikzpicture}[line width=1.1pt]
  \coordinate (B) at (0.4,0);
  \coordinate (A) at (1.5,0);
  \coordinate (O) at (2.6,0);
  \tkzDrawLine[tmblue,add=0.6 and 0.6](B,O)
  \tkzDrawPoints[size=4.5,color=tmpoint,fill=tmpoint](B,A,O)
  \tkzLabelPoint[below=3pt](B){$B$}
  \tkzLabelPoint[below=3pt](A){$A$}
  \tkzLabelPoint[below=3pt](O){$O$}
  \node[tminc] at (3.3,0.28) {$a$};
\end{tikzpicture}
\end{document}
